\chapter{总结与展望}

\section{工作总结}

本作品围绕“接收者正确、发送者也正确”这一通信需求，完成了抗量子身份基匹配加密的通用设计、
多算法实现和系统优化。发送者用身份指定接收对象，接收者在解密时指定认可的发送来源；两项
身份同时匹配，系统才恢复消息。

协议以两层匿名HIBE控制身份路径，以IBS认证发送者。32字节随机种子解决HIBE短明文与大签名、
长消息之间的容量矛盾；HKDF和AES-256-GCM承担应用载荷；一次性公钥承诺与外层签名把内部身份
认证连接到完整密文字节。建立、双方发钥、加密和匹配解密五个算法在所有组合中保持一致。

软件层将HIBE、IBS、组合配置、公共协议、状态管理和应用入口分开。两种HIBE与五种常规IBS
形成十种端到端配对，覆盖NTRU结构、普通矩阵格、环格和模格路线。目标规模Wang--PW进一步
与LATTE完成端到端组合。新增算法通过接口、规范对象和能力声明接入，不需要复制公共协议。

系统已经提供命令行、浏览器演示、固定发钥和跨进程恢复。使用者可以从算法选择开始，完成系统
建立、角色密钥生成、消息发送、身份匹配解密和错误输入测试。公开对象、角色密钥和密文均采用
带版本的规范编码，可以在不同进程间传递。

\section{创新点及实际效果}

作品的创新不只体现在接入算法数量，更体现在如何把数学方案变成可运行、可扩展的完整系统。
表~\ref{tab:innovation-summary}按照“问题—方法—结果—代价”归纳主要工作。

\begin{table}[H]
\centering\small
\caption{主要创新点、结果与使用取舍}
\label{tab:innovation-summary}
\begin{tabularx}{\textwidth}{@{}>{\raggedright\arraybackslash}p{.18\textwidth}>{\raggedright\arraybackslash}p{.29\textwidth}>{\raggedright\arraybackslash}p{.28\textwidth}>{\raggedright\arraybackslash}X@{}}
\toprule
\rowcolor{black!8} 问题 & 方法 & 结果 & 代价或适用方式 \\
\midrule
不同HIBE／IBS难以直接组合 & 独立接口、SuiteConfig、能力预检、短种子Hybrid与统一字节认证 & 十种常规配对和LATTE＋Wang--PW复用同一五算法 & 新后端需实现严格对象并声明容量 \\
Sageloli稠密展开不可承受 & 保留环多项式，逐频率执行低秩平方根与采样计划 & PARAM-I、II均完成建立、发钥、签名和验证；RSS约1.36、6.19 GiB & 需要数值变换计划和精确方程复核 \\
SRGG大整数对象存储冗长 & 按字段范围使用定长、Zigzag、二元位打包 & 主状态缩小5.05--5.46倍，签名缩小2.55--2.65倍 & 编解码增加少量处理时间 \\
Lai固定身份对反复委派 & 接收会话预先生成并验证叶密钥 & 准备约510秒后，后续解密中位数28.6秒 & 额外保存约34.7 MiB会话叶钥 \\
目标规模矩阵SIS对象巨大 & 流式对象、独立有限Gaussian、固定发行与角色导入 & Wang--PW n512完成签验、恢复和LATTE组合 & 用户钥约2.06 GiB，签名约176秒、峰值约41.3 GiB \\
\bottomrule
\end{tabularx}
\end{table}

\textbf{通用设计。} 通用性首先来自密码边界清晰。HIBE只负责两层身份路径和种子封装，IBS只
负责发送身份认证，组合层只使用公开方法和对象字节。LATTE、Lai以及六种IBS可以保留各自的
模数、维度、环和采样方式。容量预检把原语差异提前转化为可计算的应用消息上限。

\textbf{结构化计算。} Sageloli优化从环方程本身出发，不把多项式机械展开为巨型矩阵。PARAM-I
的相关扰动维数达到403,456，完整实矩阵估算约1.18~TiB；逐频率计算把核心分解降为
$65\times65$小矩阵，并在系数域继续执行精确验签关系。

\textbf{对象级优化。} SRGG编码针对四类对象逐字段确定范围。同一数学对象可以在参考JSON与紧凑
二进制之间无损往返。主状态中大量三值陷门系数只占2 bit，签名中的挑战只占1 bit，从而把存储
收益落到实际密钥文件和密文载荷。

\textbf{使用模式优化。} Lai会话区分了“固定身份对的准备”和“每条消息必须执行的认证解密”。
它没有减少每条消息的验签、GCM和字节认证，而是消除同一路径的重复委派。这一方法特别适合
机构内长期连接和固定设备对。

\section{实验结论}

A、B固定同一LATTE HIBE，因此能够直接观察IBS差异。A的重复加密与解密中位数分别为
78.5毫秒和152毫秒，适合交互演示与高频收发。B的重复加密约9.17秒，时间主要花在SRGG签名；
重复解密约161毫秒，与A接近。该结果说明，B的主要成本在发送端而不是接收端。

C更换为Lai C2与Sageloli PARAM-I。一次完整解密约425秒，其中Lai委派约393秒。接收会话把
固定路径的委派移到准备阶段后，后续解密中位数降至28.6秒。C的完整密文约4.34~MiB，主要由
Sageloli签名和Lai胶囊组成。

Wang--PW目标档展示了另一种资源分布。Setup约536.63秒，单用户发钥约5,431.95秒；复用已发
用户密钥后，一次签名约176秒，LATTE组合加密约192.54秒，匹配解密约2.85秒。发送端的签名
和用户密钥是主要开销，接收端仍能在数秒内完成公开验签与解封装。

密文组成分析还指出一个与算法选择同样重要的问题：A组96~KiB固定空间来自Lamport一次性公钥
和签名，占32字节消息密文的约69.2\%。C组则由大尺寸IBS签名和HIBE胶囊主导。后续缩短通信量
时，应根据组合选择不同目标，而不是只压缩应用消息。

\section{后续工作}

下一批实验将在Intel Core Ultra~9 275HX同一台机器、同一运行时和统一输入下，对十种常规组合
各补测10组独立状态。测量将覆盖Setup、双方发钥、首次与重复收发、恢复、对象尺寸和RSS，并把
首轮、热态和恢复后操作分开统计。完成后，实验章将用这组矩阵作为十种组合的统一横向结果，
A、B、C专项实验继续用于解释优化过程。

通信方面，A组最值得继续研究的是替换或压缩当前96~KiB的一次性认证结构，同时保持内部身份
签名与完整密文字节之间的绑定。C组可以进一步压缩Sageloli签名和Lai胶囊，并增加多套独立主
密钥下的会话测量。Wang--PW则适合继续优化用户陷门表示、流式签名准备和发钥资源调度。

应用方面，可以把当前本地发钥机构扩展为独立服务，由业务系统完成账户审核和身份字符串映射，
发送端与接收端只持有各自角色对象。现有规范编码、SuiteConfig和角色导入已经提供了这一部署
方式所需的基础。

总体而言，本作品已经从通用密码构造走到多原语、完整协议和可操作软件。不同格结构的方案能够
在同一系统中选择和比较，主要计算与存储瓶颈也有对应的实现方法和实验结果。这些工作为后续
参数完善、统一补测和面向具体场景的部署提供了稳定基础。
